\topic{本章实验、故意破坏与练习}

\begin{experimentbox}{核对一次用户 ABI 调用}
选择 \code{write}，记录系统调用号、参数寄存器、用户指针验证、返回寄存器和
errno 转换。用用户程序传入正常缓冲与跨页坏指针。
\end{experimentbox}

\begin{faultinjectionbox}{相信用户长度并直接复制}
传入跨越用户地址上限的 pointer+length。预期发生整数回绕或内核页访问。恢复
为先检查加法溢出、canonical、VMA 和逐页权限。
\end{faultinjectionbox}

\begin{pitfallbox}
ABI 固定宽度和布局；用户指针始终不可信；procfs 文本是一次快照，不应暴露
悬空内核指针。
\end{pitfallbox}

\textbf{练习}
\begin{enumerate}[leftmargin=2.2em]
  \item x86-64 syscall 返回值通常放在哪个寄存器？
  \item pointer+length 为什么先查溢出？
  \item 用户结构为何用显式宽度字段？
\end{enumerate}
\begin{answerbox}
1. RAX；2. 回绕后范围检查会错误接受小终点；3. 固定跨编译器/平台 ABI。
\end{answerbox}
